%%=============================================================================
%% Conclusie
%%=============================================================================

\chapter{Conclusie}%
\label{ch:conclusie}

% TODO: Trek een duidelijke conclusie, in de vorm van een antwoord op de
% onderzoeksvra(a)g(en). Wat was jouw bijdrage aan het onderzoeksdomein en
% hoe biedt dit meerwaarde aan het vakgebied/doelgroep? 
% Reflecteer kritisch over het resultaat. In Engelse teksten wordt deze sectie
% ``Discussion'' genoemd. Had je deze uitkomst verwacht? Zijn er zaken die nog
% niet duidelijk zijn?
% Heeft het onderzoek geleid tot nieuwe vragen die uitnodigen tot verder 
%onderzoek?

In deze bachelorstudie werd onderzocht in welke mate een XR-applicatie op de Meta Quest 3 een veilige, gebruiksvriendelijke en objectief meetbare leeromgeving kan creëren voor Basic Life Support training voor studenten verpleegkunde. Door de samenvatting van de literatuurstudie, het technische ontwerp en de daaropvolgende ontwikkeling van de Proof of Concept kunnen de onderzoeksvragen op de volgende manier worden beantwoord.

\section{Beantwoording van de onderzoeksvragen}
De technische en didactische analyse laat zien dat de Meta Quest 3, samen met een gelaagde structuur die is gebaseerd op een State Machine, in staat is om realistische medische simulaties te faciliteren die voldoen aan de onderwijsbehoeften van het verpleegkundig onderwijs. De applicatie ondersteunt niet alleen de procedurele precisie volgens de ERC-richtlijnen, maar levert door de immersieve werking van XR ook een contextuele meerwaarde die doorgaanse simulatielabs vaak missen.  \cite{ERC2025BLS, Pottle2019}.

In termen van schaalbaarheid en objectiviteit is aangetoond dat het toepassen van een gedetailleerd dataloggingsysteem, dat is gebaseerd op milliseconde-accurate timestamps en JSON-serialisatie, mogelijk maakt om kwantitatieve informatie te genereren over de uitvoering van taken. Deze gegevens voldoen aan de kwaliteitseisen voor officiële beoordelingen en vormen een basis voor de toekenning van CME-credits in de toekomst. De technische resultaten tonen aan dat een modulaire opbouw binnen de Unity-engine de benodigde flexibiliteit biedt om de PoC te verder te ontwikkelen naar een breder gebruiksvriendelijk platform binnen HOGENT. \cite{Radianti2020VRHE}.

\section{Discussie en kritische reflectie}
Tijdens het nadenken over de uitkomsten van deze bachelorproef kan worden vastgesteld dat de oorspronkelijke verwachtingen over de technische haalbaarheid zijn vervuld. De Meta Quest 3 toont aan een effectief platform te zijn voor zelfstandige medische simulatie, waarbij de handmatige tracking voldoende nauwkeurig is voor het valideren van procedures. Toch moet kritisch worden opgemerkt dat het ontbreken van fysieke tegendruk een inherent obstakel blijft voor de volledige vervanging van traditionele technieken. Hoewel de literatuur aantoont dat XR effectief is voor cognitieve en procedurele retentie, blijft de psychomotorische overgang naar de werkelijkheid zonder haptische weerstand een onderwerp van discussie. \cite{Alaker2016VR, ILCOR2023Immersive}.

Er was verwacht dat een volledig softwarematige oplossing kan voldoen aan strenge medische protocollen, maar de uitdaging om dit binnen een standalone headset te verbeteren zonder de framerate te verliezen, bleek een aanzienlijke technische uitdaging te zijn. Ondanks dat de applicatie technisch sterk is, blijft de langetermijninvloed op de leerresultaten bij een grote groep studenten in de verpleegkundeopleiding een onderwerp dat aanvullend onderzoek vereist. Dit onderzoek levert geen bijdrage door bestaande tools te vervangen, maar door een schaalbare, datagedreven aanvulling te bieden die de kans op herhaaldelijk oefenen aanzienlijk vermindert. \cite{Kyaw2019VR, Padilha2019VR}.

\section{Beperkingen van het onderzoek}
Hoewel de PoC de technische en procedurele uitvoerbaarheid bewijst, heeft dit onderzoek enkele bewuste beperkingen. In het kader van deze bachelorproef is de beoordeling beperkt tot conceptuele validatie, expertbeoordelingen en scenario-walkthroughs, waarbij grootschalige gebruikerstests met studenten als essentieel vervolgonderzoek worden gezien. Bovendien blijft de biomechanische uitvoering van compressies een abstractie in de virtuele wereld, waardoor de applicatie op dit moment fungeert als een didactisch hulpmiddel voor procedurele training en niet als een uiteindelijke test voor fysieke vaardigheden. \cite{ILCOR2023Immersive}.

\section{Aanbevelingen en verder onderzoek}
Het onderzoek heeft resulteerd in nieuwe vragen die aanmoedigen tot verdere studie. Een belangrijke vraag is op welke manier de combinatie van hybride simulatie, waarbij de XR-omgeving wordt verbonden met fysieke manikins, de leerervaring beïnvloedt in vergelijking met een enkel virtuele omgeving. Voor de verdere groei binnen de beoogde 'XR Academy' is het raadzaam om zich te richten op multi-user scenario's, waarbij communicatie en teamwerk de belangrijkste rol spelen. \cite{Liaw2018MUVW}.
 
Bovendien is een klinische validatie vereist, waarbij de invloed op de leerretentie en simulatorziekte bij studenten op een gestructureerde manier wordt gemeten met behulp van de SUS- en SSQ-vragenlijsten \cite{SUSBrooke1996, Kennedy1993SSQ}. XR SkillSim vormt een solide fundament voor een hedendaagse benadering van medische vaardigheidstraining, die de verbinding legt tussen theoretische kennis en de ingewikkelde werkelijkheid van de klinische praktijk.

